\chapter{Grass Allergy}

\section*{Definition and Overview}
Grass allergy is a type of allergic rhinitis (hay fever) caused by an immune reaction to pollen released by grasses. It is one of the most common pollen allergies in Australia, and it typically causes symptoms during the spring and summer grass pollen season, although in some regions grass pollen can cause problems for much of the year. Symptoms include sneezing, a runny or blocked nose, itchy and watery eyes, and sometimes asthma symptoms such as coughing and wheezing. Grass allergy is managed with allergen avoidance, medicines such as antihistamines and nasal sprays, and in selected people, immunotherapy (allergy vaccination) that can reduce sensitivity long term.

\section*{Epidemiology}
Grass pollen allergy affects a substantial proportion of Australians. Around 1 in 5 Australians has allergic rhinitis, and grass pollen is a leading trigger. The most important grass species in Australia include ryegrass, which is widespread in temperate regions, and subtropical grasses such as Bahia and Johnson grass in northern areas. Grass pollen is a major trigger of thunderstorm asthma, a phenomenon in which storms rupture pollen grains and cause severe asthma epidemics, as occurred in Melbourne in 2016. Pollen counts are highest in spring and early summer, and symptoms are worst on windy, dry days.

\section*{Aetiology and Pathophysiology}
Grass allergy is an immediate-type (IgE-mediated) allergic reaction to proteins in grass pollen.

\begin{itemize}
    \item \textbf{Pollen:} grasses release large amounts of pollen, which is carried by the wind over long distances
    \item \textbf{Sensitisation:} in susceptible people, the immune system produces antibodies (IgE) against grass pollen proteins
    \item \textbf{Allergic response:} on re-exposure, pollen triggers mast cells to release histamine and other chemicals, causing sneezing, itching, and mucus production
    \item \textbf{Eye involvement:} pollen contacts the eyes, causing allergic conjunctivitis
    \item \textbf{Airway involvement:} in people with asthma, pollen can trigger bronchospasm, and in thunderstorm asthma, ruptured pollen grains penetrate deep into the lungs
    \item \textbf{Cross-reactivity:} some people with grass allergy also react to foods such as tomato, potato, and melon (pollen-food syndrome), causing an itchy mouth and throat
\end{itemize}

\section*{Clinical Features}
Symptoms follow the grass pollen season and are often predictable.

\begin{itemize}
    \item Sneezing, often in bouts
    \item A runny, itchy, or blocked nose
    \item Itchy, red, watery, or puffy eyes
    \item Itching of the roof of the mouth, throat, or ears
    \item Post-nasal drip and throat clearing
    \item Cough, wheezing, and chest tightness in people with asthma
    \item Fatigue and poor concentration from disturbed sleep
    \item Symptoms worse outdoors, on windy days, and in the morning or evening
\end{itemize}

\section*{Differential Diagnosis}
Other causes of similar symptoms should be considered.

\begin{itemize}
    \item Other pollen allergies, including tree and weed pollen
    \item Dust mite and animal dander allergies, which tend to be year-round
    \item Mould allergy
    \item Non-allergic (vasomotor) rhinitis, triggered by temperature change, fumes, or stress
    \item The common cold and other viral infections
    \item Sinusitis, which can follow rhinitis
    \item Medication-induced rhinitis
    \item Structural nasal problems, including a deviated septum and nasal polyps
\end{itemize}

\section*{When to Seek Medical Care}
See a doctor if symptoms are frequent, severe, or affect sleep, work, or school, or if over-the-counter medicines are not controlling them. People with asthma should ensure their asthma is well managed during the pollen season.

\textbf{Call 000 (or local emergency services) for:}
\begin{itemize}
    \item Severe wheezing, breathlessness, or difficulty breathing during the pollen season, which may indicate thunderstorm asthma or a severe asthma attack
    \item Asthma that is not responding to your reliever inhaler
    \item Signs of anaphylaxis, including difficulty breathing, swelling of the face or throat, or collapse, though this is rare with grass pollen alone
\end{itemize}

\section*{Diagnosis and Investigations}
Grass allergy is confirmed by correlating symptoms with the pollen season and by allergy testing.

\begin{itemize}
    \item \textbf{Clinical history:} seasonal symptoms and response to treatment are highly suggestive
    \item \textbf{Skin prick testing:} small drops of grass pollen extract are placed on the skin and pricked; a raised red bump indicates allergy
    \item \textbf{Blood tests:} specific IgE (allergy antibody) levels to grass pollen can be measured
    \item \textbf{Pollen counts:} checking daily pollen forecasts helps confirm the link and plan avoidance
    \item \textbf{Lung function tests:} if asthma is suspected
    \item \textbf{Review by an allergist:} for complex or severe cases
\end{itemize}

\section*{Management}
Management combines avoidance, medicines, and immunotherapy.

\begin{itemize}
    \item \textbf{Avoidance:} limit outdoor activity when pollen counts are high, keep windows closed, and shower and change clothes after being outdoors
    \item \textbf{Antihistamines:} non-drowsy oral antihistamines control sneezing, itching, and runny nose
    \item \textbf{Nasal sprays:} corticosteroid nasal sprays are the most effective treatment for nasal symptoms and are best used regularly through the season
    \item \textbf{Eye drops:} antihistamine eye drops relieve itchy, watery eyes
    \item \textbf{Combined treatment:} many people need both an antihistamine and a nasal spray
    \item \textbf{Asthma care:} people with asthma should take their preventer medicines and have an asthma action plan
    \item \textbf{Immunotherapy:} allergy vaccination (tablets or injections) over several years can reduce sensitivity to grass pollen and is effective for many people
    \item \textbf{Saline rinses:} washing the nose with salt water clears pollen and soothes the lining
\end{itemize}

\section*{Complications}
\begin{itemize}
    \item Reduced quality of life, sleep disturbance, and poor concentration
    \item Sinusitis from prolonged nasal blockage
    \item Worsening of asthma, including severe attacks during thunderstorm asthma
    \item Ear infections and eustachian tube dysfunction in children
    \item Side effects of medicines, including drowsiness with older antihistamines
    \item Effects on work and school performance during the season
\end{itemize}

\section*{Prognosis and Outcomes}
Grass allergy is a chronic condition, but it is highly manageable. Most people control symptoms well with regular use of appropriate medicines through the pollen season. Immunotherapy offers the possibility of long-term improvement, with many people experiencing substantially reduced symptoms and medicine use for years after completing treatment. The most important precautions are good asthma control and awareness of high-risk days, because the main serious complication, thunderstorm asthma, is largely preventable with timely reliever and preventer use.

\section*{Prevention and Self-Care}
\begin{itemize}
    \item Check daily pollen counts and limit outdoor time on high-count days
    \item Keep windows and doors closed during the season, and use air conditioning if available
    \item Wear sunglasses to protect the eyes from pollen
    \item Shower and change clothes after being outdoors, and wash hair before bed
    \item Dry bedding indoors rather than on an outdoor line
    \item Start nasal sprays before the season begins and use them regularly
    \item Take antihistamines as directed and choose non-drowsy formulations
    \item If you have asthma, take your preventer daily and keep your reliever with you
    \item Have an asthma action plan and know what to do in a thunderstorm asthma event
    \item Ask your doctor whether immunotherapy is suitable for you
\end{itemize}

\section*{Sources and Further Reading}
The following reputable sources provide further information for healthcare students and patients seeking reliable, current advice about grass allergy.

\begin{itemize}
    \item Australasian Society of Clinical Immunology and Allergy. \textit{Hay fever (allergic rhinitis) information.} https://www.allergy.org.au/
    \item Melbourne Pollen Count (University of Melbourne). \textit{Pollen forecasts and thunderstorm asthma.} https://www.melbournepollen.com.au/
    \item National Asthma Council Australia. \textit{Thunderstorm asthma and pollen information.} https://www.nationalasthma.org.au/
    \item Healthdirect Australia. \textit{Hay fever and pollen allergy.} https://www.healthdirect.gov.au/
    \item NHS. \textit{Hay fever: overview and treatment.} https://www.nhs.uk/
    \item American Academy of Allergy, Asthma and Immunology. \textit{Pollen allergy.} https://www.aaaai.org/
\end{itemize}
